\section{The Gambit Backend}
\label{sec:gambit}

The Gambit backend turns the EventGraph into an extensive-form game
in Gambit's \TT{.efg} format~\citep{McKelveyMcLennanTurocy2006Gambit}.
Once written, the file can be loaded directly into Gambit (or
OpenSpiel~\citep{LanctotEtAl2019OpenSpiel}, via standard adapters)
to compute Nash equilibria, dominant strategies, or any other
property Gambit knows how to check.  This section describes how the
backend constructs the tree from the graph.

\subsection{Frontier exploration}

The backend is a depth-first frontier expander.  At a configuration
$C$, it enumerates the frontier, groups its nodes by owning role,
and for each role that has at least one node on the frontier it
generates a Gambit decision node.  Multiple roles may have nodes on
the same frontier; these are simultaneous moves in the strategic
analysis (\S\ref{sec:eventgraph:frontier}).  Gambit's \TT{.efg}
format does not have a primitive simultaneous block: the backend
nests the per-role decision nodes in a canonical role order and
arranges the information sets so that each role's information set
hides the other roles' choices at the same strategic move.  The
strategic content is in the information sets, not the textual
nesting order.

For each role, the actions are the products of values its nodes may
take.  Bounded-type fields have an explicit finite alphabet
(e.g.\ \TT{door = \{0,1,2\}} has three values); the \TT{Quit} move
is added as an extra alphabet symbol when the role's node has a
quit handler.  A node performing a hidden commit has its commit
parameter listed as ``Hidden(\textit{value})'' rather than the
cleartext value: from the perspective of the game, the value
matters but the observer does not see it.  When the corresponding
reveal node fires, the resolved value collapses ``Hidden(\dots)''
to a concrete number.

For each combination of role moves at the current frontier, the
backend computes the resulting configuration, recurses, and emits
the subtree.  When the frontier is empty --- the configuration has
finished every node --- the \TT{withdraw} clause is evaluated to
produce a terminal node carrying the payoff vector.

\subsection{Information sets}

A player's information set at a decision node is the set of
histories they cannot distinguish, given what they have observed.
The backend computes this directly from
$\mathrm{obs}_R(B)$ (\S\ref{sec:eventgraph:visibility-recall}):
two Gambit decision nodes are in the same information set for
role $R$ iff the projections of their pre-state to
$R$-observable fields are equal.

In practice, hidden commits create the information sets that
matter.  In the Vickrey example, each bidder's reveal node is at a
depth where the other bidders' commitments have fired but their
values have not been revealed; the reveal node therefore sits in
an information set whose members differ only in the other
bidders' hidden values.  This is the correct strategic
representation of a sealed-bid auction.

\subsection{Pruning}

Two reductions are run on the produced tree before emission.

\emph{Unreachable subtrees are dropped.}  When a guard rules out a
move (the guard evaluates to false for the would-be combination of
values), no edge is emitted; this is straightforward as long as
the guard's free variables are all evaluated at the decision
node's information set.

\emph{Trivial continuations are collapsed.}  When a player has
exactly one legal move at a decision node (often because
\TT{Quit} is the only choice on a degenerate branch), the
backend emits a single edge and proceeds.  This keeps the tree
size manageable on examples with many forced timeouts.

\subsection{Limits}

A few characteristics of the backend should be acknowledged.

\paragraph{Tree size blows up with the type universe.}  The
backend enumerates the cross-product of move alphabets at each
frontier.  For a hypothetical three-bidder auction with 11-value
bids and a quit option, the commit frontier alone yields
$12^3 = 1728$ branches (where the strategic choice actually
happens); each reveal node only branches by 2 (the committed
value or Quit), but the reveal frontier is visited once per
commit combination, so the total tree size grows
multiplicatively.  The bundled Vickrey example uses a 3-value
bid type to keep the tree analysable.  This is a limit of
tree-enumeration analysis, not of the language: a Vegas
program with a larger bid type compiles and deploys correctly,
but the resulting EFG becomes too large for Gambit-style
equilibrium search, and the analyst gains nothing from the EFG
backend past that point.  This is fine for the example library;
it is not what one would deploy against a competitive-bidding
mechanism with millions of bid levels.  Vegas's contract-side bid types are typically narrow because
Gambit cannot otherwise analyse them; wider ranges deploy fine
but make Gambit-side analysis infeasible.  This is the explicit
tradeoff in carrying both views in the same source.

\paragraph{Move alphabets must be bounded.}  An action whose
parameter has type \TT{int} is rejected by the Gambit backend at
compile time.  The Solidity backend has no such restriction (it
stores an \TT{int256} directly); the Gambit constraint is
enforced via a per-backend disable list in the test harness.

\paragraph{Quit appears as a real move.}  Because the bail-out
subgame is a strategic option in Vegas's semantics, the
\TT{.efg} contains a \TT{Quit} branch wherever the source action
has an or-handler.  This is the right thing
game-theoretically (it captures the strategic option of timing
out), but it does inflate the tree.

\paragraph{Chance nodes.}  Actions owned by a chance label ---
a \TT{random} role or a \TT{sample} binding --- are emitted as
Gambit chance moves.  Each parameter carries a per-node
distribution: either the explicit \TT{$\sim$ uniform/weighted}
annotation declared at the parameter, or the inferred uniform
over the parameter's bounded domain when no annotation is
present.  A self-only guard on the chance action is interpreted
as standard probabilistic conditioning: the guard-surviving
support is renormalised so that probabilities sum to one on the
children of the chance node.  Contextual guards (referring to
prior fields) are conditioned per trace, with a fresh
renormalisation at each reached context.  Each chance event is
a node of its own, with the same simultaneous-move treatment as a
player move.
